\documentclass[../../main.tex]{subfiles}
\graphicspath{{\subfix{../../image/}}}

\begin{document}

\section{临界点}\label{section:ch08-5-临界点}

至此我们研究了在可能受约束条件下寻找各种能量泛函极小元的问题，并发现它们所满足的相应非线性 Euler--Lagrange 方程。本节转向寻找 Euler--Lagrange PDE 的其他解的问题：通过寻找其他临界点。这些临界点一般而言不是极小元，而是 $I[\cdot]$ 的"鞍点"。

\subsection{山路定理}\label{subsection:ch08-5-1-山路定理}

接着开发一些工具，以保证一个抽象泛函 $I[\cdot]$ 存在临界点。

\subsubsection*{a. 临界点，形变}

今后 $H$ 表示实 Hilbert 空间，范数 $\|\cdot\|$，内积 $(,)$。设 $I:H\to\mathbb{R}$ 为 $H$ 上的非线性泛函。

\begin{definition}\label{definition:ch08-5-可微}
称 $I$ 在 $u\in H$ 处可微，若存在 $v\in H$ 使
\begin{align}
I[w]=I[u]+(v,w-u)+o(\|w-u\|)\qquad (w\in H).\label{align:ch08-5-eq01}
\end{align}
元素 $v$（若存在）唯一。于是写 $I'[u]=v$。
\end{definition}

\begin{definition}\label{definition:ch08-5-C1}
称 $I\in C^1(H;\mathbb{R})$，若对每个 $u\in H$，$I'[u]$ 存在，且映射 $I':H\to H$ 连续。
\end{definition}

下面发展的理论在 $I\in C^1(H;\mathbb{R})$ 时成立，但若此外假设
\begin{align}
I':H\to H\text{ 在 }H\text{ 的有界子集上 Lipschitz 连续},\label{align:ch08-5-eq02}
\end{align}
证明将大为简化。

\paragraph{记号。}
(i) 记 $\mathcal{C}$ 为满足\eqref{align:ch08-5-eq02} 的函数 $I\in C^1(H;\mathbb{R})$ 之集。
(ii) 若 $c\in\mathbb{R}$，写
\begin{align*}
A_c\triangleq\{u\in H\mid I[u]\leqslant c\},\qquad K_c\triangleq\{u\in H\mid I[u]=c,\ I'[u]=0\}.
\end{align*}

\paragraph{定义。}
(i) 称 $u\in H$ 为\textbf{临界点}，若 $I'[u]=0$。
(ii) 实数 $c$ 是\textbf{临界值}，若 $K_c\neq\emptyset$。

现要证明：若 $c$ 不是临界值，则可把集合 $A_{c+\varepsilon}$ 优美地形变为 $A_{c-\varepsilon}$（某 $\varepsilon>0$）。想法是在 $H$ 中解一个适当的 ODE 并沿所得流"下山"。由于 $H$ 一般无穷维，需要某种紧性条件。

\begin{definition}\label{definition:ch08-5-PalaisSmale}
泛函 $I\in C^1(H;\mathbb{R})$ 满足\textbf{Palais--Smale 紧性条件}，若每个满足
(i) $\{I[u_k]\}_{k=1}^\infty$ 有界，且
(ii) $I'[u_k]\to 0$（在 $H$ 中）
的序列 $\{u_k\}_{k=1}^\infty\subset H$ 在 $H$ 中预紧。
\end{definition}

\begin{theorem}[形变定理]\label{theorem:ch08-5-1-形变定理}
设 $I\in\mathcal{C}$ 满足 Palais--Smale 条件。又设
\begin{align}
K_c=\emptyset.\label{align:ch08-5-eq03}
\end{align}
则对每个充分小的 $\varepsilon>0$，存在常数 $0<\delta<\varepsilon$ 及函数
\begin{align*}
\boldsymbol{\eta}\in C([0,1]\times H;H)
\end{align*}
使映射
\begin{align*}
\boldsymbol{\eta}_t(u)=\boldsymbol{\eta}(t,u)\qquad (0\leqslant t\leqslant 1,\ u\in H)
\end{align*}
满足
(i) $\boldsymbol{\eta}_0(u)=u$（$u\in H$），
(ii) $\boldsymbol{\eta}_1(u)=u$（$u\notin I^{-1}[c-\varepsilon,c+\varepsilon]$），
(iii) $I[\boldsymbol{\eta}_t(u)]\leqslant I[u]$（$u\in H$，$0\leqslant t\leqslant 1$），
(iv) $\boldsymbol{\eta}_1(A_{c+\delta})\subset A_{c-\delta}$。
\begin{proof}
1. 首先断言存在常数 $0<\sigma,\varepsilon<1$ 使
\begin{align}
\|I'[u]\|\geqslant\sigma\quad\text{对每个 }u\in A_{c+\varepsilon}\setminus A_{c-\varepsilon}.\label{align:ch08-5-eq04}
\end{align}
反证法。若\eqref{align:ch08-5-eq04} 对所有常数 $\sigma,\varepsilon>0$ 不成立，则存在序列 $\sigma_k\to 0$、$\varepsilon_k\to 0$ 及元素
\begin{align}
u_k\in A_{c+\varepsilon_k}\setminus A_{c-\varepsilon_k}\label{align:ch08-5-eq05}
\end{align}
满足
\begin{align}
\|I'[u_k]\|\leqslant\sigma_k.\label{align:ch08-5-eq06}
\end{align}
由 Palais--Smale 条件，存在子列 $\{u_{k_j}\}_{j=1}^\infty$ 及元素 $u\in H$ 使 $u_{k_j}\to u$ 在 $H$ 中。但由于 $I\in C^1(H;\mathbb{R})$，\eqref{align:ch08-5-eq05} 与\eqref{align:ch08-5-eq06} 蕴涵 $I[u]=c$，$I'[u]=0$。故 $K_c\neq\emptyset$，与假设\eqref{align:ch08-5-eq03} 矛盾。
2. 现固定 $\delta$ 满足
\begin{align}
0<\delta<\varepsilon,\qquad 0<\delta<\frac{\sigma^2}{2}.\label{align:ch08-5-eq07}
\end{align}
写
\begin{align*}
A\triangleq\{u\in H\mid I[u]\leqslant c-\varepsilon\text{ 或 }I[u]\geqslant c+\varepsilon\},\\
B\triangleq\{u\in H\mid c-\delta\leqslant I[u]\leqslant c+\delta\}.
\end{align*}
由于 $I'$ 在有界集上有界，可验证映射 $u\mapsto\operatorname{dist}(u,A)+\operatorname{dist}(u,B)$ 在 $H$ 的每个有界子集下方有界于某正常数。因此函数
\begin{align*}
g(u)\triangleq\frac{\operatorname{dist}(u,A)}{\operatorname{dist}(u,A)+\operatorname{dist}(u,B)}\qquad (u\in H)
\end{align*}
在有界集上 Lipschitz 连续，且满足
\begin{align}
0\leqslant g\leqslant 1,\quad g=0\text{ 在 }A\text{ 上},\quad g=1\text{ 在 }B\text{ 上}.\label{align:ch08-5-eq08}
\end{align}
令
\begin{align}
h(t)\triangleq\begin{cases}1,&0\leqslant t\leqslant 1,\\1/t,&t\geqslant 1.\end{cases}\label{align:ch08-5-eq09}
\end{align}
最后定义映射 $V:H\to H$ 为
\begin{align}
V(u)\triangleq -g(u)h(\|I'[u]\|)I'[u]\qquad (u\in H).\label{align:ch08-5-eq10}
\end{align}
注意 $V$ 有界。
3. 现对每个 $u\in H$ 考虑 ODE
\begin{align}
\begin{cases}
\frac{\mathrm{d}\boldsymbol{\eta}}{\mathrm{d}t}(t)=V(\boldsymbol{\eta}(t))&(t>0),\\
\boldsymbol{\eta}(0)=u.
\end{cases}\label{align:ch08-5-eq11}
\end{align}
由于 $V$ 有界且在有界集上 Lipschitz 连续，存在唯一解，对所有 $t\geqslant 0$ 存在。写 $\boldsymbol{\eta}=\boldsymbol{\eta}(t,u)=\boldsymbol{\eta}_t(u)$（$t\geqslant 0$，$u\in H$）以显式显示解对时间 $t$ 与初位置 $u\in H$ 的依赖。限于时间 $0\leqslant t\leqslant 1$，如此定义的映射 $\boldsymbol{\eta}\in C([0,1]\times H;H)$ 满足断言 (i)、(ii)。
4. 现计算
\begin{align}
\frac{\mathrm{d}}{\mathrm{d}t}I[\boldsymbol{\eta}_t(u)]&=(I'[\boldsymbol{\eta}_t(u)],\frac{\mathrm{d}}{\mathrm{d}t}\boldsymbol{\eta}_t(u))\notag\\
&=(I'[\boldsymbol{\eta}_t(u)],V(\boldsymbol{\eta}_t(u)))\notag\\
&=-g(\boldsymbol{\eta}_t(u))h(\|I'[\boldsymbol{\eta}_t(u)]\|)\|I'[\boldsymbol{\eta}_t(u)]\|^2.\label{align:ch08-5-eq12}
\end{align}
特别地
\begin{align*}
\frac{\mathrm{d}}{\mathrm{d}t}I[\boldsymbol{\eta}_t(u)]\leqslant 0\qquad (u\in H,\ 0\leqslant t\leqslant 1),
\end{align*}
故断言 (iii) 成立。
5. 现固定任一点
\begin{align}
u\in A_{c+\delta}.\label{align:ch08-5-eq13}
\end{align}
要证
\begin{align}
\boldsymbol{\eta}_1(u)\in A_{c-\delta}\label{align:ch08-5-eq14}
\end{align}
从而验证断言 (iv)。若对某 $0\leqslant t\leqslant 1$，$\boldsymbol{\eta}_t(u)\notin B$，则已完成；故不妨假设 $\boldsymbol{\eta}_t(u)\in B$（$0\leqslant t\leqslant 1$）。于是 $g(\boldsymbol{\eta}_t(u))=1$（$0\leqslant t\leqslant 1$）。由计算\eqref{align:ch08-5-eq12} 得
\begin{align}
\frac{\mathrm{d}}{\mathrm{d}t}I[\boldsymbol{\eta}_t(u)]=-h(\|I'[\boldsymbol{\eta}_t(u)]\|)\|I'[\boldsymbol{\eta}_t(u)]\|^2.\label{align:ch08-5-eq15}
\end{align}
若 $\|I'[\boldsymbol{\eta}_t(u)]\|\geqslant 1$，则\eqref{align:ch08-5-eq09} 与\eqref{align:ch08-5-eq04} 蕴涵
\begin{align*}
\frac{\mathrm{d}}{\mathrm{d}t}I[\boldsymbol{\eta}_t(u)]=-\|I'[\boldsymbol{\eta}_t(u)]\|\leqslant-\sigma<-\sigma^2.
\end{align*}
反之，若 $\|I'[\boldsymbol{\eta}_t(u)]\|\leqslant 1$，\eqref{align:ch08-5-eq09} 与\eqref{align:ch08-5-eq04} 给出
\begin{align*}
\frac{\mathrm{d}}{\mathrm{d}t}I[\boldsymbol{\eta}_t(u)]\leqslant-\sigma^2.
\end{align*}
这两个不等式及\eqref{align:ch08-5-eq15} 然后蕴涵
\begin{align*}
I[\boldsymbol{\eta}_1(u)]\leqslant I[u]-\sigma^2\leqslant c+\delta-\sigma^2\leqslant c-\delta\quad\text{由 }\eqref{align:ch08-5-eq07}.
\end{align*}
此估计建立\eqref{align:ch08-5-eq14} 并完成证明。
\end{proof}
\end{theorem}

\subsubsection*{b. 山路定理}

接着利用一个有趣的"min-max"技巧，用上面构造的形变 $\boldsymbol{\eta}$ 推得临界点的存在性。

\begin{theorem}[山路定理]\label{theorem:ch08-5-2-山路定理}
设 $I\in\mathcal{C}$ 满足 Palais--Smale 条件。又设
(i) $I[0]=0$，
(ii) 存在常数 $r,a>0$ 使
\begin{align*}
I[u]\geqslant a\quad\text{若 }\|u\|=r,
\end{align*}
且
(iii) 存在元素 $v\in H$ 满足
\begin{align*}
\|v\|>r,\qquad I[v]\leqslant 0.
\end{align*}
定义
\begin{align*}
\Gamma\triangleq\{\mathbf{g}\in C([0,1];H)\mid \mathbf{g}(0)=0,\ \mathbf{g}(1)=v\}.
\end{align*}
则
\begin{align*}
c=\inf_{\mathbf{g}\in\Gamma}\max_{0\leqslant t\leqslant 1}I[\mathbf{g}(t)]
\end{align*}
是 $I$ 的临界值。

把 $I[\cdot]$ 的图像想象成一片风景：在 $0$ 处有一个低谷，周围环绕一圈山脉；山脉之外在 $v$ 处另有一个低谷。想法是寻找连接 $0$ 与 $v$ 的路径 $\mathbf{g}$，它经过一个山口，即 $I[\cdot]$ 的一个鞍点。但请注意：我们只断言在"能量水平" $c$ 处存在一个临界点，它未必对应真正的鞍点。
\begin{proof}
显然
\begin{align}
c\geqslant a.\label{align:ch08-5-eq16}
\end{align}
设 $c$ 不是 $I$ 的临界值，即
\begin{align}
K_c=\emptyset.\label{align:ch08-5-eq17}
\end{align}
然后取任一充分小的数
\begin{align}
0<\varepsilon<\frac a2.\label{align:ch08-5-eq18}
\end{align}
由形变\cref{theorem:ch08-5-1-形变定理}，存在常数 $0<\delta<\varepsilon$ 及同胚 $\boldsymbol{\eta}:H\to H$ 满足
\begin{align}
\boldsymbol{\eta}(A_{c+\delta})\subset A_{c-\delta}\label{align:ch08-5-eq19}
\end{align}
且
\begin{align}
\boldsymbol{\eta}(u)=u\quad\text{若 }u\notin I^{-1}[c-\varepsilon,c+\varepsilon].\label{align:ch08-5-eq20}
\end{align}
现选 $\mathbf{g}\in\Gamma$ 满足
\begin{align}
\max_{0\leqslant t\leqslant 1}I[\mathbf{g}(t)]\leqslant c+\delta.\label{align:ch08-5-eq21}
\end{align}
则 $\widehat{\mathbf{g}}\triangleq\boldsymbol{\eta}\circ\mathbf{g}$ 也属于 $\Gamma$，因为由\eqref{align:ch08-5-eq20}，$\boldsymbol{\eta}(\mathbf{g}(0))=\boldsymbol{\eta}(0)=0$ 且 $\boldsymbol{\eta}(\mathbf{g}(1))=\boldsymbol{\eta}(v)=v$。但随后\eqref{align:ch08-5-eq21} 蕴涵 $\max_{0\leqslant t\leqslant 1}I[\widehat{\mathbf{g}}(t)]\leqslant c-\delta$，故
\begin{align*}
c=\inf_{\mathbf{g}\in\Gamma}\max_{0\leqslant t\leqslant 1}I[\mathbf{g}(t)]\leqslant c-\delta,
\end{align*}
矛盾。
\end{proof}
\end{theorem}

\subsection{对半线性椭圆 PDE 的应用}\label{subsection:ch08-5-2-半线性椭圆}

为说明山路定理的用处，现考察半线性边值问题：
\begin{align}
\begin{cases}
-\Delta u=f(u)&\text{在 }U\text{ 中},\\
u=0&\text{在 }\partial U\text{ 上}.
\end{cases}\label{align:ch08-5-eq22}
\end{align}
设 $f$ 光滑，且对某
\begin{align*}
1<p<\frac{n+2}{n-2}
\end{align*}
有
\begin{align}
|f(z)|\leqslant C(1+|z|^p),\qquad |f'(z)|\leqslant C(1+|z|^{p-1})\qquad (z\in\mathbb{R}),\label{align:ch08-5-eq23}
\end{align}
其中 $C$ 为常数。还假设
\begin{align}
0\leqslant F(z)\leqslant\gamma f(z)z\quad\text{对某常数 }\gamma<\frac12,\label{align:ch08-5-eq24}
\end{align}
其中 $F(z)\triangleq\int_0^z f(s)\,\mathrm{d}s$，$z\in\mathbb{R}$。最后对常数 $0<a\leqslant A$ 假设
\begin{align}
a|z|^{p+1}\leqslant|F(z)|\leqslant A|z|^{p+1}\qquad (z\in\mathbb{R}).\label{align:ch08-5-eq25}
\end{align}

现\eqref{align:ch08-5-eq25} 蕴涵 $f(0)=0$，故显然 $u\equiv 0$ 是\eqref{align:ch08-5-eq22} 的平凡解。我们要找另一个。

注意到 PDE
\begin{align*}
-\Delta u=|u|^{p-1}u
\end{align*}
满足上述假设。我们将在 \S9.4.2 再次回到这一特殊非线性。

\begin{theorem}[存在性]\label{theorem:ch08-5-3-半线性存在}
边值问题\eqref{align:ch08-5-eq22} 至少有一个弱解 $u\not\equiv 0$。
\begin{proof}
1. 定义
\begin{align}
I[u]\triangleq\int_U\frac12|Du|^2-F(u)\,\mathrm{d}x\label{align:ch08-5-eq26}
\end{align}
（$u\in H_0^1(U)$）。我们要对 $I[\cdot]$ 应用山路定理。设 $H=H_0^1(U)$，范数 $\|u\|=\left(\int_U|Du|^2\,\mathrm{d}x\right)^{1/2}$，内积 $(u,v)=\int_U Du\cdot Dv\,\mathrm{d}x$。则
\begin{align*}
I[u]=\frac12\|u\|^2-\int_U F(u)\,\mathrm{d}x\triangleq I_1[u]-I_2[u].
\end{align*}
2. 首先断言
\begin{align}
I\text{ 属于类 }\mathcal{C}.\label{align:ch08-5-eq27}
\end{align}
为此注意对每个 $u,w\in H$，
\begin{align*}
I_1[w]=\frac12\|w\|^2=\frac12\|u+w-u\|^2=\frac12\|u\|^2+(u,w-u)+\frac12\|w-u\|^2.
\end{align*}
故 $I_1$ 在 $u$ 处可微，$I_1'[u]=u$。因此 $I_1\in\mathcal{C}$。
3. 接着考察项 $I_2$。回顾 Lax--Milgram 定理（\S6.2.1）：对每个 $v^*\in H^{-1}(U)$，问题
\begin{align*}
\begin{cases}
-\Delta v=v^*&\text{在 }U\text{ 中},\\
v=0&\text{在 }\partial U\text{ 上}
\end{cases}
\end{align*}
有唯一解 $v\in H_0^1(U)$。写 $v=Kv^*$，于是
\begin{align}
K:H^{-1}(U)\to H_0^1(U)\text{ 是等距同构。}\label{align:ch08-5-eq28}
\end{align}
特别注意若 $w\in L^{\frac{2n}{n+2}}(U)$，则由
\begin{align*}
\langle w^*,u\rangle\triangleq\int_U wu\,\mathrm{d}x\qquad (u\in H_0^1(U))
\end{align*}
定义的线性泛函 $w^*$ 属于 $H^{-1}(U)$。（我们将滥用记号说"$w\in H^{-1}(U)$"。）接着注意 $p\cdot\frac{2n}{n+2}<\frac{n+2}{n-2}\cdot\frac{2n}{n+2}=2^*$，故若 $u\in H_0^1(U)$，则 $f(u)\in L^{\frac{2n}{n+2}}(U)\subset H^{-1}(U)$。
现证明若 $u\in H_0^1(U)$，则
\begin{align}
I_2'[u]=K[f(u)].\label{align:ch08-5-eq29}
\end{align}
为此首先注意
\begin{align*}
F(a+b)=F(a)+f(a)b+\int_0^1(1-s)f'(a+sb)\,\mathrm{d}s\,b^2.
\end{align*}
于是对每个 $w\in H_0^1(U)$，
\begin{align}
I_2[w]&=\int_U F(w)\,\mathrm{d}x=\int_U F(u+w-u)\,\mathrm{d}x\notag\\
&=\int_U F(u)+f(u)(w-u)\,\mathrm{d}x+R\notag\\
&=I_2(u)+\int_U DK[f(u)]\cdot D(w-u)\,\mathrm{d}x+R,\label{align:ch08-5-eq30}
\end{align}
其中由\eqref{align:ch08-5-eq23}，余项 $R$ 满足
\begin{align*}
|R|&\leqslant C\int_U(1+|u|^{p-1}+|w-u|^{p-1})|w-u|^2\,\mathrm{d}x\\
&\leqslant C\left(\int_U|w-u|^2+|w-u|^{p+1}\,\mathrm{d}x\right)\\
&\quad+C\left(\int_U|u|^{p+1}\,\mathrm{d}x\right)^{\frac{p-1}{p+1}}\left(\int_U|w-u|^{p+1}\,\mathrm{d}x\right)^{\frac{2}{p+1}}.
\end{align*}
由于 $p+1<2^*$，Sobolev 不等式给出 $R=o(\|w-u\|)$。故由\eqref{align:ch08-5-eq28} 见
\begin{align*}
I_2[w]=I_2[u]+(K[f(u)],w-u)+o(\|w-u\|),
\end{align*}
即所证。
最后注意若 $u,\overline{u}\in H_0^1(U)$ 且 $\|u\|,\|\overline{u}\|\leqslant L$，则
\begin{align*}
\|I_2'[u]-I_2'[\overline{u}]\|&=\|K[f(u)]-K[f(\overline{u})]\|_{H_0^1(U)}\\
&=\|f(u)-f(\overline{u})\|_{H^{-1}(U)}\\
&\leqslant\|f(u)-f(\overline{u})\|_{L^{\frac{2n}{n+2}}(U)}.
\end{align*}
但
\begin{align*}
\|f(u)-f(\overline{u})\|_{L^{\frac{2n}{n+2}}(U)}&\leqslant C\left(\int_U\big((1+|u|^{p-1}+|\overline{u}|^{p-1})|u-\overline{u}|\big)^{\frac{2n}{n+2}}\,\mathrm{d}x\right)^{\frac{n+2}{2n}}\\
&\leqslant C\left(\int_U(1+|u|^{p-1}+|\overline{u}|^{p-1})^{\frac{n+2}{4}}\,\mathrm{d}x\right)^{\frac2n}\|u-\overline{u}\|_{L^{2^*}(U)}\\
&\leqslant C(L)\|u-\overline{u}\|_{L^{2^*}(U)}\leqslant C(L)\|u-\overline{u}\|,
\end{align*}
其中用了\eqref{align:ch08-5-eq23}。故 $I_2':H_0^1(U)\to H_0^1(U)$ 在有界集上 Lipschitz 连续。因此 $I_2\in\mathcal{C}$，断言\eqref{align:ch08-5-eq27} 得证。
4. 现验证 Palais--Smale 条件。为此设 $\{u_k\}_{k=1}^\infty\subset H_0^1(U)$，且
\begin{align}
\{I[u_k]\}_{k=1}^\infty\text{ 有界}\label{align:ch08-5-eq31}
\end{align}
且
\begin{align}
I'[u_k]\to 0\quad\text{在 }H_0^1(U)\text{ 中}.\label{align:ch08-5-eq32}
\end{align}
由上述
\begin{align}
u_k-K(f(u_k))\to 0\quad\text{在 }H_0^1(U)\text{ 中}.\label{align:ch08-5-eq33}
\end{align}
于是对每个 $\varepsilon>0$，当 $k$ 充分大时
\begin{align*}
|(I'[u_k],v)|=\left|\int_U Du_k\cdot Dv-f(u_k)v\,\mathrm{d}x\right|\leqslant\varepsilon\|v\|\qquad (v\in H_0^1(U)).
\end{align*}
在上式中令 $v=u_k$，得
\begin{align*}
\left|\int_U|Du_k|^2-f(u_k)u_k\,\mathrm{d}x\right|\leqslant\varepsilon\|u_k\|
\end{align*}
对每个 $\varepsilon>0$ 及所有充分大的 $k$ 成立。特别取 $\varepsilon=1$，见
\begin{align}
\int_U f(u_k)u_k\,\mathrm{d}x\leqslant\|u_k\|^2+\|u_k\|\label{align:ch08-5-eq34}
\end{align}
对所有充分大的 $k$ 成立。但由于\eqref{align:ch08-5-eq31} 说
\begin{align*}
\left(\frac12\|u_k\|^2-\int_U F(u_k)\,\mathrm{d}x\right)\leqslant C<\infty
\end{align*}
对所有 $k$ 及某常数 $C$，推得
\begin{align*}
\|u_k\|^2\leqslant C+2\int_U F(u_k)\,\mathrm{d}x\leqslant C+2\gamma(\|u_k\|^2+\|u_k\|)
\end{align*}
此由\eqref{align:ch08-5-eq34}、\eqref{align:ch08-5-eq24}。
由于 $2\gamma<1$，发现 $\{u_k\}_{k=1}^\infty$ 在 $H_0^1(U)$ 中有界。故存在子列 $\{u_{k_j}\}_{j=1}^\infty$ 及 $u\in H_0^1(U)$，使 $u_{k_j}\rightharpoonup u$ 弱收敛于 $H_0^1(U)$ 且 $u_{k_j}\to u$ 在 $L^{p+1}(U)$ 中（后一断言成立是因为 $p+1<2^*$）。但随后 $f(u_k)\to f(u)$ 在 $H^{-1}(U)$ 中，故 $K[f(u_k)]\to K[f(u)]$ 在 $H_0^1(U)$ 中。因此\eqref{align:ch08-5-eq33} 蕴涵
\begin{align}
u_{k_j}\to u\quad\text{在 }H_0^1(U)\text{ 中}.\label{align:ch08-5-eq35}
\end{align}
5. 最后验证山路定理其余假设。显然 $I[0]=0$。现设 $u\in H_0^1(U)$，$\|u\|=r$，其中 $r>0$ 待下方选取。则
\begin{align}
I[u]=I_1[u]-I_2[u]=\frac{r^2}{2}-I_2[u].\label{align:ch08-5-eq36}
\end{align}
假设\eqref{align:ch08-5-eq25} 蕴涵，由于 $p+1<2^*$，
\begin{align*}
|I_2[u]|\leqslant C\int_U|u|^{p+1}\,\mathrm{d}x\leqslant C\left(\int_U|u|^{2^*}\,\mathrm{d}x\right)^{\frac{p+1}{2^*}}\leqslant C\|u\|^{p+1}\leqslant Cr^{p+1}.
\end{align*}
于是由\eqref{align:ch08-5-eq36}，
\begin{align*}
I[u]\geqslant\frac{r^2}{2}-Cr^{p+1}\geqslant\frac{r^2}{4}=a>0,
\end{align*}
只要 $r>0$ 充分小，因为 $p+1>2$。现固定某 $u\in H$，$u\not\equiv 0$。写 $v=tu$，$t>0$ 待选。则
\begin{align*}
I[v]&=I_1[tu]-I_2[tu]=t^2I_1[u]-\int_U F(tu)\,\mathrm{d}x\\
&\leqslant t^2I_1[u]-at^{p+1}\int_U|u|^{p+1}\,\mathrm{d}x\quad\text{由 }\eqref{align:ch08-5-eq25}\\
&<0
\end{align*}
对 $t>0$ 充分大。
6. 终于检验了山路定理的全部假设。故必存在函数 $u\in H_0^1(U)$，$u\not\equiv 0$，满足
\begin{align*}
I'[u]=u-K[f(u)]=0.
\end{align*}
特别对每个 $v\in H_0^1(U)$，有
\begin{align*}
\int_U Du\cdot Dv\,\mathrm{d}x=\int_U f(u)v\,\mathrm{d}x,
\end{align*}
故 $u$ 是\eqref{align:ch08-5-eq22} 的弱解。
\end{proof}
\end{theorem}

进一步关于非线性 Poisson 方程的讨论、特别是假设\eqref{align:ch08-5-eq23} 中临界指数 $\frac{n+2}{n-2}$ 的意义，见 \S9.4.2。

\end{document}
